\documentclass[12pt,a4paper]{article}

\usepackage[utf8]{inputenc}
\usepackage[vietnamese]{babel}
\usepackage{amsmath,amssymb,amsthm}
\usepackage{geometry}
\usepackage{enumitem}
\usepackage[most]{tcolorbox}
\usepackage{hyperref}
\usepackage{fancyhdr}
\usepackage{graphicx}
\usepackage{tikz}
\usepackage{eso-pic}
\usepackage{xcolor}

\geometry{a4paper, margin=2cm, bottom=2.5cm}
\hypersetup{colorlinks=true, linkcolor=blue!70!black, urlcolor=blue}
\setlist[itemize]{topsep=4pt, itemsep=2pt}
\setlist[enumerate]{topsep=4pt, itemsep=2pt}

\AddToShipoutPictureFG{%
  \AtPageCenter{%
    \begin{tikzpicture}[remember picture, overlay]
      \node[rotate=45, scale=1, opacity=0.06]{\fontsize{60}{72}\selectfont\bfseries\sffamily Đặng Tấn Quí};
    \end{tikzpicture}%
  }%
}

\pagestyle{fancy}
\fancyhf{}
\fancyhead[L]{\small\textit{Đại số cơ sở}}
\fancyhead[R]{\small\textit{Đặng Tấn Quí}}
\fancyfoot[C]{\thepage}
\fancyfoot[L]{\footnotesize\textcopyright\ Đặng Tấn Quí}
\fancyfoot[R]{\footnotesize SĐT: 0377\,122\,210}
\renewcommand{\headrulewidth}{0.4pt}
\renewcommand{\footrulewidth}{0.4pt}

\fancypagestyle{plain}{\fancyhf{}\fancyfoot[C]{\thepage}\fancyfoot[L]{\footnotesize\textcopyright\ Đặng Tấn Quí}\fancyfoot[R]{\footnotesize SĐT: 0377\,122\,210}\renewcommand{\headrulewidth}{0pt}\renewcommand{\footrulewidth}{0.4pt}}

\newtcolorbox{dinhnghia}[1][]{enhanced, breakable, colback=blue!3!white, colframe=blue!60!black, fonttitle=\bfseries, title={Định nghĩa}, #1}
\newtcolorbox{dinhly}[1][]{enhanced, breakable, colback=green!3!white, colframe=green!50!black, fonttitle=\bfseries, title={Định lý}, #1}
\newtcolorbox{tinhchat}[1][]{enhanced, breakable, colback=yellow!5!white, colframe=orange!50!black, fonttitle=\bfseries, title={Tính chất}, #1}
\newtcolorbox{phuongphap}[1][]{enhanced, breakable, colback=gray!5!white, colframe=gray!50!black, fonttitle=\bfseries, title={Phương pháp}, #1}
\newtcolorbox{luuy}[1][]{enhanced, breakable, colback=red!3!white, colframe=red!50!black, fonttitle=\bfseries, title={Lưu ý}, #1}
\newtcolorbox{congthuc}[1][]{enhanced, breakable, colback=cyan!3!white, colframe=cyan!50!black, fonttitle=\bfseries, title={Công thức}, #1}
\newtcolorbox{vidu}[1][]{enhanced, breakable, colback=violet!3!white, colframe=violet!50!black, fonttitle=\bfseries, title={Ví dụ}, #1}

\begin{document}

\thispagestyle{empty}
\begin{tikzpicture}[remember picture, overlay]
  \fill[blue!80!black] (current page.south west) rectangle (current page.north east);
  \node[anchor=west, white, font=\fontsize{26}{32}\bfseries\selectfont]
    at ([xshift=3cm, yshift=-7.5cm]current page.north west) {ĐẠI SỐ CƠ SỞ};
  \node[anchor=west, white, font=\fontsize{18}{24}\selectfont]
    at ([xshift=3cm, yshift=-9cm]current page.north west) {Tập hợp, ánh xạ, số phức, cấu trúc đại số --- Năm 1, HK1};
  \node[anchor=west, white, font=\Large\bfseries] at ([xshift=3cm, yshift=-13cm]current page.north west) {Biên soạn:};
  \node[anchor=west, yellow!80!white, font=\fontsize{22}{28}\bfseries\selectfont] at ([xshift=3cm, yshift=-14.5cm]current page.north west) {ĐẶNG TẤN QUÍ};
  \node[anchor=south east, white, font=\Large\bfseries] at ([xshift=-2cm, yshift=3cm]current page.south east) {\the\year};
\end{tikzpicture}

\newpage
\tableofcontents
\newpage

%% ================================================================
%%  CHƯƠNG 1: MỆNH ĐỀ VÀ LOGIC
%% ================================================================
\section{Mệnh đề và logic}

\subsection{Mệnh đề và giá trị chân lý}

\begin{dinhnghia}
  \textbf{Mệnh đề} là một khẳng định chỉ có thể đúng hoặc sai.
  
  Giá trị chân lý: \textbf{đúng} (1, true) hoặc \textbf{sai} (0, false). 
  
  Mệnh đề không thể vừa đúng vừa sai.
\end{dinhnghia}

\begin{vidu}
  ``$3 + 5 = 8$'' (đúng)
  
  ``$2 > 5$'' (sai)
  
  ``Số 7 là số nguyên tố'' (đúng).

  ``$2x = 6$ đúng khi $x = 4$'' (sai).
  
  Không phải mệnh đề: ``$x + 1 = 0$'' (chứa biến, chưa xác định).
  
  Không phải mệnh đề: ``$x^3 - 4x^2 + 5x - 6$'' (chứa biến, chưa xác định).
\end{vidu}

\subsection{Các phép toán logic}

\begin{dinhnghia}
  Cho hai mệnh đề $A$, $B$:
  \begin{itemize}
    \item \textbf{Phủ định} $\neg A$: đúng khi $A$ sai, sai khi $A$ đúng.
    \item \textbf{Hội} (và) $A \wedge B$: đúng khi cả $A$ và $B$ đều đúng.
    \item \textbf{Tuyển} (hoặc) $A \vee B$: sai khi cả $A$ và $B$ đều sai.
    \item \textbf{Kéo theo} $A \Rightarrow B$: sai khi $A$ đúng và $B$ sai; đúng trong các trường hợp còn lại.
    \begin{itemize}
        \item A là điều kiện đủ để có B
        \item B là điều kiện cần để có A
    \end{itemize}
    \item \textbf{Tương đương} $A \Leftrightarrow B$: đúng khi $A$ và $B$ cùng đúng hoặc cùng sai.
    \begin{itemize}
        \item A là điều kiện cần và đủ để có B
        \item B cũng là điều kiện cần và đủ để có A
    \end{itemize}
  \end{itemize}
\end{dinhnghia}

\begin{congthuc}[title={Bảng chân lý}]
  \begin{center}
  \renewcommand{\arraystretch}{1.3}
  \begin{tabular}{|c|c||c|c|c|c|c|}
  \hline
  $A$ & $B$ & $\neg A$ & $A \wedge B$ & $A \vee B$ & $A \Rightarrow B$ & $A \Leftrightarrow B$ \\
  \hline
  1 & 1 & 0 & 1 & 1 & 1 & 1 \\
  1 & 0 & 0 & 0 & 1 & 0 & 0 \\
  0 & 1 & 1 & 0 & 1 & 1 & 0 \\
  0 & 0 & 1 & 0 & 0 & 1 & 1 \\
  \hline
  \end{tabular}
  \end{center}
\end{congthuc}

\begin{luuy}
  $A \Rightarrow B$ chỉ sai khi ``giả thiết đúng nhưng kết luận sai''. 
  
  Khi $A$ sai thì $A \Rightarrow B$ luôn đúng (vô nghĩa nhưng quy ước).
\end{luuy}

\subsection{Tính chất các phép logic}

\begin{tinhchat}
  \begin{itemize}
    \item Phủ định của phủ định: $\neg(\neg A) \Leftrightarrow A$.
    \item Giao hoán: $A \wedge B \Leftrightarrow B \wedge A$.
    \item Kết hợp: $(A \wedge B) \wedge C \Leftrightarrow A \wedge (B \wedge C)$.
    \item Phân phối: $A \vee (B \wedge C) \Leftrightarrow (A \vee B) \wedge (A \vee C)$, $\quad A \wedge (B \vee C) \Leftrightarrow (A \wedge B) \vee (A \wedge C)$.
    \item Biểu diễn kéo theo: $A \Rightarrow B \Leftrightarrow (\neg A) \vee B \Leftrightarrow \neg(A \wedge (\neg B))$.
    \item \textbf{Luật De Morgan}: $\neg(A \wedge B) \Leftrightarrow (\neg A) \vee (\neg B)$, $\quad \neg(A \vee B) \Leftrightarrow (\neg A) \wedge (\neg B)$.
    \item \textbf{Nguyên lý phản đảo}: $A \Rightarrow B \Leftrightarrow (\neg B) \Rightarrow (\neg A)$.
  \end{itemize}
\end{tinhchat}

\begin{vidu}
  Chứng minh ``Nếu $n^2$ chẵn thì $n$ chẵn'' tương đương với ``Nếu $n$ lẻ thì $n^2$ lẻ'' (phản đảo dễ chứng minh hơn).
\end{vidu}

\subsection{Mệnh đề với lượng từ}

\begin{dinhnghia}
  \textbf{Hàm mệnh đề} $P(x)$: phát biểu phụ thuộc biến $x$, trở thành mệnh đề khi thay $x$ bằng giá trị cụ thể.
\end{dinhnghia}

\begin{dinhnghia}
  \begin{itemize}
    \item \textbf{Mệnh đề toàn thể}: ``$\forall x \in X$: $P(x)$'' (với mọi $x$, $P(x)$ đúng).
    \item \textbf{Mệnh đề tồn tại}: ``$\exists x \in X$: $P(x)$'' (tồn tại ít nhất một $x$ sao cho $P(x)$ đúng).
  \end{itemize}
\end{dinhnghia}

\begin{congthuc}[title={Phủ định}]
  $\neg(\forall x: P(x)) \Leftrightarrow \exists x: \neg P(x)$.
  
  $\neg(\exists x: P(x)) \Leftrightarrow \forall x: \neg P(x)$.
\end{congthuc}

\begin{vidu}
  Phủ định của ``Mọi số thực đều có bình phương không âm'' là ``Tồn tại số thực có bình phương âm''. Phủ định của ``Tồn tại $x$ sao cho $x^2 = -1$'' là ``Với mọi $x$, $x^2 \ne -1$''.
\end{vidu}

%% ================================================================
%%  CHƯƠNG 2: TẬP HỢP
%% ================================================================
\section{Tập hợp}

\subsection{Các khái niệm cơ bản}

\begin{dinhnghia}
  \textbf{Tập hợp} là một khái niệm cơ bản của toán học, là một nhóm các đối tượng, các đối tượng gọi là \textbf{các phần tử} của tập hợp.
  \begin{itemize}
    \item $x \in A$: $x$ thuộc $A$ (hoặc $x$ là phần tử của $A$).
    \item $x \notin A$: $x$ không thuộc $A$.
  \end{itemize}
\end{dinhnghia}

\begin{dinhnghia}
  \textbf{Tập con}: $A \subset B$ (hoặc $A \subseteq B$) $\Leftrightarrow$ $\forall x \in A \Rightarrow x \in B$.
  
  \textbf{Bằng nhau}: $A = B$ $\Leftrightarrow$ $A \subset B$ và $B \subset A$.

  \textbf{Tập con thật sự}: $A \subset B$ và $A \ne B$

  Không phải tập con: $A \not \subset B$
\end{dinhnghia}

\begin{dinhnghia}
  \textbf{Tập rỗng} $\varnothing$: tập không chứa phần tử nào. Với mọi tập $A$: $\varnothing \subset A$.
\end{dinhnghia}

\subsection{Cách xác định tập hợp}

\begin{phuongphap}[title={Liệt kê}]
  Liệt kê tất cả các phần tử trong dấu ngoặc nhọn. Ví dụ: $A = \{1, 2, 3, 4\}$, $B = \{a, b, c\}$.
\end{phuongphap}

\begin{phuongphap}[title={Chỉ ra tính chất đặc trưng}]
  $A := \{x \mid x \text{ có tính chất } P\}$ hoặc $A := \{x : P(x)\}$.
  
  Ví dụ: $E := \{x \in \mathbb{N} \mid x \text{ chẵn và } x < 10\} = \{0, 2, 4, 6, 8\}$.
\end{phuongphap}

\begin{luuy}
  Ký hiệu $:=$ nghĩa là ``được định nghĩa bởi''.
  
  Có thể biểu diễn tập hợp qua \textbf{biểu đồ Ven} (hình tròn minh họa giao, hợp, phần bù).
\end{luuy}

\begin{vidu}
  \begin{itemize}
    \item $\mathbb{N} = \{0, 1, 2, 3, \ldots\}$: tập số tự nhiên.
    \item $\mathbb{N}^* = \{1, 2, 3, \ldots\}$: tập số tự nhiên khác $0$.
    \item $\mathbb{Z} = \{\ldots, -2, -1, 0, 1, 2, \ldots\}$: tập số nguyên.
    \item $\mathbb{Q} := \{\dfrac{p}{q} \mid q \ne 0,\; p \in \mathbb{Z},\; q \in \mathbb{Z}\}$: tập số hữu tỉ.
    \item $\mathbb{R}$: tập số thực.
  \end{itemize}
\end{vidu}

\subsection{Các phép toán trên tập hợp}

\begin{dinhnghia}
  Cho $A$, $B$ là các tập con của tập $E$ (không gian):
  \begin{itemize}
    \item \textbf{Hợp}: $A \cup B = \{x \mid x \in A \text{ hoặc } x \in B\}$.
    \item \textbf{Giao}: $A \cap B = \{x \mid x \in A \text{ và } x \in B\}$.
    \begin{itemize}
        \item A và B rời nhau khi $A \cap B = \varnothing$
    \end{itemize}
    \item \textbf{Hiệu}: $A \setminus B = \{x \in A \mid x \notin B\}$.
    \item \textbf{Phần bù} (trong $E$): $\overline{A} = C_E A = E \setminus A$ (khi $A \subset E$).
  \end{itemize}
\end{dinhnghia}

\begin{tinhchat}[title={Tính chất các phép toán}]
  \begin{itemize}
    \item Luỹ đẳng: $A \cup A = A$, $A \cap A = A$
    \item Giao hoán: $A \cup B = B \cup A$, $A \cap B = B \cap A$.
    \item Kết hợp: $(A \cup B) \cup C = A \cup (B \cup C)$, $(A \cap B) \cap C = A \cap (B \cap C)$.
    \item Phân phối: $A \cup (B \cap C) = (A \cup B) \cap (A \cup C)$, $A \cap (B \cup C) = (A \cap B) \cup (A \cap C)$.
    \item Luật De Morgan: $\overline{A \cup B} = \overline{A} \cap \overline{B}$, $\overline{A \cap B} = \overline{A} \cup \overline{B}$.
  \end{itemize}
\end{tinhchat}

\begin{phuongphap}[title={Chứng minh De Morgan}]
  Chứng minh $\overline{A \cup B} = \overline{A} \cap \overline{B}$:
  \[
    x \in \overline{A \cup B} \;\Leftrightarrow\; x \notin A \cup B \;\Leftrightarrow\; x \notin A \text{ và } x \notin B \;\Leftrightarrow\; x \in \overline{A} \text{ và } x \in \overline{B} \;\Leftrightarrow\; x \in \overline{A} \cap \overline{B}.
  \]

  Chứng minh $\overline{A \cap B} = \overline{A} \cup \overline{B}$:
  \[
    x \in \overline{A \cap B} \;\Leftrightarrow\; x \notin A \cap B \;\Leftrightarrow\; x \notin A \text{ hoặc } x \notin B \;\Leftrightarrow\; x \in \overline{A} \text{ hoặc } x \in \overline{B} \;\Leftrightarrow\; x \in \overline{A} \cup \overline{B}.
  \]
\end{phuongphap}

\begin{luuy}
    \begin{itemize}
        \item $(A \setminus B) \cap (B \setminus A) = \varnothing$
    \end{itemize}
\end{luuy}

\subsection{Hợp và giao của họ tập hợp}

\begin{dinhnghia}
  Cho họ tập hợp $\{A_i\}_{i \in I}$ (chỉ số trong tập $I \subset \mathbb{N}$):
  \begin{itemize}
    \item \textbf{Hợp}: $\displaystyle\bigcup_{i \in I} A_i = \{x \mid \exists i \in I: x \in A_i\}$.
    \item \textbf{Giao}: $\displaystyle\bigcap_{i \in I} A_i = \{x \mid \forall i \in I: x \in A_i\}$.
  \end{itemize}
  Khi $I = \{1, 2, \ldots, n\}$ viết $\bigcup_{i=1}^{n} A_i$, $\bigcap_{i=1}^{n} A_i$.
\end{dinhnghia}

\begin{tinhchat}[title={Luật De Morgan cho họ tập hợp}]
  Trong không gian $E$: $\displaystyle\overline{\bigcup_{i \in I} A_i} = \bigcap_{i \in I} \overline{A_i}$, $\quad \displaystyle\overline{\bigcap_{i \in I} A_i} = \bigcup_{i \in I} \overline{A_i}$.
\end{tinhchat}

\begin{vidu}
  $A_1 = [0,1]$, $A_2 = [1,2]$, $A_3 = [2,3]$. Khi đó $\bigcup_{i=1}^{3} A_i = [0,3]$. $\bigcap_{i=1}^{3} A_i = \varnothing$ vì không có số thực nào thuộc cả ba đoạn.
\end{vidu}

\subsection{Phủ và Phân hoạch}

\begin{dinhnghia}
  Họ các tập con $\{A_1, A_2, \ldots, A_n\}$ của $X$ gọi là \textbf{phủ} của $X$ nếu:
  \begin{itemize}
    \item $\bigcup_{i=1}^{n} A_i = X$ (hợp bằng toàn bộ $X$).
  \end{itemize}
  
  Họ các tập con $\{A_1, A_2, \ldots, A_n\}$ của $X$ gọi là \textbf{phân hoạch} của $X$ nếu:
  \begin{itemize}
    \item Họ các tập con đó là \textbf{phủ} của $X$.
    \item Và $A_i \cap A_j = \varnothing$ với mọi $i \ne j$ (đôi một rời nhau).
  \end{itemize}
\end{dinhnghia}

\begin{vidu}
  Phân hoạch $\mathbb{Z}$: $A_0 = \{\ldots, -4, -2, 0, 2, 4, \ldots\}$ (chẵn), $A_1 = \{\ldots, -3, -1, 1, 3, \ldots\}$ (lẻ). Mỗi số nguyên thuộc đúng một trong hai tập.
\end{vidu}

\subsection{Tích Đề-các}

\begin{dinhnghia}
  \textbf{Tích Đề-các} (tích Descartes) của $A$ và $B$ (theo thứ tự ấy):
  \[
    A \times B := \{(a, b) \mid a \in A,\; b \in B\}.
  \]
  Phần tử $(a, b)$ gọi là \textbf{cặp thứ tự}: $(a, b) = (c, d)$ $\Leftrightarrow$ $a = c$ và $b = d$.
\end{dinhnghia}

\begin{dinhnghia}
  \textbf{Tích Đề-các của $n$ tập hợp}: $A_1 \times A_2 \times \cdots \times A_n = \{(a_1, a_2, \ldots, a_n) \mid a_i \in A_i,\; i = 1, \ldots, n\}$. Ký hiệu $A^n = A \times A \times \cdots \times A$ ($n$ lần).
\end{dinhnghia}

\begin{vidu}
  $\mathbb{R} \times \mathbb{R} = \mathbb{R}^2$ (mặt phẳng). $\mathbb{R}^3$ (không gian ba chiều). 
  
  $\{0,1\}^n$: tập các dãy nhị phân độ dài $n$.
\end{vidu}

\begin{luuy}
  Nói chung $A \times B \ne B \times A$. 
  
  Ví dụ: $\{1\} \times \{2\} = \{(1,2)\}$ còn $\{2\} \times \{1\} = \{(2,1)\}$.
\end{luuy}

%% ================================================================
%%  CHƯƠNG 3: QUAN HỆ HAI NGÔI
%% ================================================================
\section{Quan hệ hai ngôi}

\begin{dinhnghia}
  \textbf{Quan hệ hai ngôi}: Cho tập $E$ có tính chất $\mathcal{R}$ liên quan đến 2 phần tử của E. Nếu a và b là 2 phần tử của E thỏa mãn tính chất $\mathcal{R}$ thì ta nói a quan hệ $\mathcal{R}$ với b viết $a\mathcal{R}b$.

  Cặp $(a, b)$ là một tập con của $E \times E$ thỏa mãn quan hệ $\mathcal{R}$.

  Đồ thị quan hệ $\mathcal{R}$ là tập hợp các cặp $(a, b)$ đó.
\end{dinhnghia}

\begin{vidu}
  Trên tập các đường thẳng trong không gian. "Đường thẳng $(d)$ vuông góc với đường thẳng $(d')$" là 1 quan hệ 2 ngôi.

  Đồ thị quan hệ "$a = b$" trên $\mathbb{R}$ là đường phân giác của góc vuông $I$ và $III$ trong mặt phẳng tọa độ $Oab$
\end{vidu}

\begin{dinhnghia}[title={Các tính chất}]
  Cho quan hệ $\mathcal{R}$ trên $X$:
  \begin{itemize}
    \item \textbf{Phản xạ}: $\forall x \in X$: $x \mathcal{R} x$.
    \item \textbf{Đối xứng}: $\forall x, y \in X$: $x \mathcal{R} y \Rightarrow y \mathcal{R} x$.
    \item \textbf{Phản đối xứng}: $\forall x, y \in X$: $x \mathcal{R} y$ và $y \mathcal{R} x \Rightarrow x = y$.
    \item \textbf{Bắc cầu}: $\forall x, y, z \in X$: $x \mathcal{R} y$ và $y \mathcal{R} z \Rightarrow x \mathcal{R} z$.
  \end{itemize}
\end{dinhnghia}

\begin{vidu}
  Trên $\mathbb{Z}$: 
  \begin{itemize}
  \item ``$a \mid b$'' (chia hết): phản xạ, phản đối xứng, bắc cầu. Không đối xứng, vì chưa chắc $b \mid a$
  \item ``$a - b$ chẵn'': phản xạ, đối xứng, bắc cầu. Không phản đối xứng, vì $a - b$ chẵn và $b - a$ chẵn thì $a = b$ chưa chắc đúng (ví dụ $a = 2$, $b = 4$).
  \item $a > b$: bắc cầu. Không phản xạ ($a > a$ sai), không đối xứng, không phản đối xứng ($a > b$ và $b > a$ không thể cùng đúng).
  \end{itemize}
\end{vidu}

\subsection{Quan hệ tương đương}

\begin{dinhnghia}
  Quan hệ $\mathcal{R}$ trên $X$ gọi là \textbf{quan hệ tương đương} nếu có tính phản xạ, đối xứng và bắc cầu. Ký hiệu: $a \sim b \;(\mathcal{R})$ hoặc $a \sim b$.
\end{dinhnghia}

\begin{dinhnghia}
  \textbf{Lớp tương đương} của $x$: $\mathcal{T}(x, \mathcal{R}) = \{y \in X \mid y \sim x \;(\mathcal{R})\}$ là gom tất cả những cái “giống x” (theo quan hệ $\mathcal{R}$) lại thành 1 nhóm.

  Quan hệ tương đương chia tập thành các nhóm không giao nhau. Tập thương $\dfrac{X}{\sim} = \{\mathcal{T}(x, \mathcal{R}) \mid x \in X\}$ là tập gồm tất cả các nhóm đó.

  $x$ là phần tử đại diện của $\mathcal{T}(x, \mathcal{R})$
\end{dinhnghia}

\begin{vidu}
    Chia 3 dư giống nhau $X = \mathbb{Z}, a \sim b \iff a \equiv b \ (\text{mod } 3)$
    
    Các lớp tương đương:
    \begin{itemize}
        \item $\mathcal{T}(0) = \{..., -3, 0, 3, 6, ...\}$
        \item $\mathcal{T}(1) = \{..., -2, 1, 4, 7, ...\}$
        \item $\mathcal{T}(2) = \{..., -1, 2, 5, 8, ...\}$
    \end{itemize}

    Vậy: $\dfrac{\mathbb{Z}}{\sim} = \{\mathcal{T}(0), \mathcal{T}(1), \mathcal{T}(2)\}$
\end{vidu}

\subsection{Quan hệ thứ tự}

\begin{dinhnghia}
  Quan hệ $\mathcal{R}$ trên $X$ gọi là \textbf{quan hệ thứ tự} (thứ tự bộ phận) nếu có tính phản xạ, phản đối xứng và bắc cầu. Khi đó $(X, \mathcal{R})$ gọi là \textbf{tập sắp thứ tự}.
\end{dinhnghia}

\begin{dinhnghia}
  Thứ tự \textbf{toàn phần} (tuyến tính): Quan hệ thứ tự $\mathcal{R}$ trên $X$ sao cho $\forall a, b \in X$: $a \mathcal{R} b$ \textbf{hoặc} $b \mathcal{R} a$ (mọi cặp phần tử đều so sánh được).
\end{dinhnghia}

\begin{vidu}
  $(\mathbb{N}, \le)$, $(\mathbb{R}, \le)$: thứ tự toàn phần. 
  
  $(\mathcal{P}(X), \subset)$: thứ tự bộ phận (không toàn phần nếu $X$ có hơn 1 phần tử: $\{1\}$ và $\{2\}$ không so sánh được).
\end{vidu}

\begin{dinhnghia}
    Tập có thứ tự: Xét tập $X$ trong đó có 1 quan hệ thứ tự $\mathcal{R}$. Cho 2 phần tử bất kỳ $a, b \in X$ nếu $a \mathcal{R} b$ ta nói $a$ có thể so sánh được với $b$. Các phần tử có thể so sánh được với nhau, sắp xếp theo 1 thứ tự xác định theo quan hệ $\mathcal{R}$.

    Nếu $\mathcal{R}$ là 1 quan hệ thứ tự toàn phần thì tất cả các phần tử của $X$ đều được sắp xếp thứ tự theo quan hệ $\mathcal{R}$. Ta nói $X$ được sắp thứ tự toàn phần.

    Nếu $\mathcal{R}$ là 1 quan hệ thứ tự không toàn phần thì chỉ có 1 số phần tử của $X$ được sắp xếp theo quan hệ $\mathcal{R}$. Ta nói $X$ là tập có thứ tự bộ phận hay $X$ được sắp thứ tự bộ phận.
\end{dinhnghia}

%% ================================================================
%%  CHƯƠNG 4: ÁNH XẠ
%% ================================================================
\section{Ánh xạ}

\subsection{Định nghĩa và các khái niệm}

\begin{dinhnghia}
  \textbf{Ánh xạ} $f: A \to B$ là quy tắc cho tương ứng mỗi $x \in A$ với đúng một phần tử $f(x) \in B$.
  \begin{itemize}
    \item $A$: \textbf{miền xác định} (tập nguồn).
    \item $B$: \textbf{tập đích}.
    \item $f(x)$: \textbf{ảnh} của $x$. Hay $x$ là \textbf{nghịch ảnh} của $f(x)$.
  \end{itemize}
\end{dinhnghia}

\begin{dinhnghia}
  \textbf{Miền giá trị} (ảnh của $f$): $\mathrm{Im}(f) = f(A) = \{f(x) \mid x \in A\} \subset B$.
\end{dinhnghia}

\begin{dinhnghia}
  Cho $f: A \to B$, $C \subset A$, $D \subset B$:
  \begin{itemize}
    \item \textbf{Ảnh} của $C$: $f(C) = \{f(x) \mid x \in C\}$.
    \item \textbf{Nghịch ảnh} của $D$: $f^{-1}(D) = \{x \in A \mid f(x) \in D\}$.
  \end{itemize}
\end{dinhnghia}

\begin{tinhchat}[title={Tính chất ảnh và nghịch ảnh}]
  Cho $f: X \to Y$, $A, B \subset X$, $C, D \subset Y$:
  \begin{itemize}
    \item $A \subset B \Rightarrow f(A) \subset f(B)$; \quad $C \subset D \Rightarrow f^{-1}(C) \subset f^{-1}(D)$.
    \item $f(A \cup B) = f(A) \cup f(B)$; \quad $f(A \cap B) \subset f(A) \cap f(B)$ (đẳng thức khi $f$ đơn ánh).
    \item $f^{-1}(C \cup D) = f^{-1}(C) \cup f^{-1}(D)$; \quad $f^{-1}(C \cap D) = f^{-1}(C) \cap f^{-1}(D)$.
    \item $f^{-1}(Y \setminus C) = X \setminus f^{-1}(C)$.
  \end{itemize}
\end{tinhchat}

\begin{vidu}
  $f: \mathbb{R} \to \mathbb{R}$, $f(x) = x^2$. 
  \begin{itemize}
      \item $f(\{-1, 1\}) = \{1\}$.
      \item $f^{-1}(\{4\}) = \{-2, 2\}$.
      \item $f^{-1}([-1, 4]) = [-2, 2]$.
  \end{itemize}
\end{vidu}

\subsection{Đơn ánh, toàn ánh, song ánh}

\begin{dinhnghia}
  Cho $f: A \to B$:
  \begin{itemize}
    \item $f$ \textbf{đơn ánh} (injective): $f(x_1) = f(x_2) \Rightarrow x_1 = x_2$. (Hai phần tử khác nhau có ảnh khác nhau.)
    \item $f$ \textbf{toàn ánh} (surjective): $\mathrm{Im}(f) = B$, tức $\forall y \in B$, $\exists x \in A$: $f(x) = y$.
    \item $f$ \textbf{song ánh} (bijective): vừa đơn ánh vừa toàn ánh.
  \end{itemize}
\end{dinhnghia}

\begin{tinhchat}
  Nếu $f: A \to B$ đơn ánh thì $f: A \to \mathrm{Im}(f)$ là toàn ánh. 
  
  Mọi song ánh đều có thể xem như đơn ánh từ $A$ lên $B$.

  Nếu $f$ song ánh thì $f^{-1}$ cũng song ánh.
\end{tinhchat}

\begin{luuy}
  Nếu phương trình $f(x) = y$ với ẩn $x$ không thể có quá 1 nghiệm với mọi $y \in B$ thì $f$ đơn ánh
  
  Nếu phương trình $f(x) = y$ với ẩn $x$ có nghiệm với mọi $y \in B$ thì $f$ toàn ánh
  
  Nếu phương trình $f(x) = y$ với ẩn $x$ có đúng 1 nghiệm với mọi $y \in B$ thì $f$ song ánh
\end{luuy}

\begin{dinhnghia}
  \textbf{Ánh xạ hằng}: $f: A \to B$ gọi là hằng nếu $\exists c \in B$: $f(x) = c$ với mọi $x \in A$. Ký hiệu $f \equiv c$.
\end{dinhnghia}

\begin{dinhnghia}
  \textbf{Ánh xạ đồng nhất} trên $A$: $\mathrm{id}_A: A \to A$, $\mathrm{id}_A(x) = x$ với mọi $x \in A$.
\end{dinhnghia}

\begin{dinhnghia}
  \textbf{Ánh xạ đặc trưng} (hàm chỉ tiêu): Cho $A \subset X$, ta có $\chi_A: X \to \{0, 1\}$: $\chi_A(x) = 1$ nếu $x \in A$, $\chi_A(x) = 0$ nếu $x \notin A$.
\end{dinhnghia}

\begin{dinhnghia}
  \textbf{Phép chiếu} thứ $i$: Cho $A_1 \times A_2 \times \cdots \times A_n$, ta có $\pi_i: A_1 \times \cdots \times A_n \to A_i$, $\pi_i(a_1, \ldots, a_n) = a_i$.
\end{dinhnghia}

\begin{dinhnghia}
  Hai ánh xạ $f, g: A \to B$ \textbf{bằng nhau} ($f = g$) nếu $f(x) = g(x)$ với mọi $x \in A$.
\end{dinhnghia}

\begin{vidu}
  $f: \mathbb{R} \to \mathbb{R}$, $f(x) = 5$ là ánh xạ hằng. 
  
  $\pi_1: \mathbb{R}^2 \to \mathbb{R}$, $\pi_1(x,y) = x$ là phép chiếu lên trục hoành.
\end{vidu}

\subsection{Ánh xạ ngược và hợp ánh xạ}

\begin{dinhnghia}
  Nếu $f: A \to B$ song ánh thì tồn tại duy nhất \textbf{ánh xạ ngược} $f^{-1}: B \to A$ sao cho $f^{-1}(f(x)) = x$ và $f(f^{-1}(y)) = y$ với mọi $x \in A$, $y \in B$.
\end{dinhnghia}

\begin{dinhnghia}
  Cho 3 tập hợp E, F, G và 2 ánh xạ: $f: E \to F$ và $g: F \to G$, ta có mỗi $x \in E$ thông qua trung gian $y \in F$ tồn tại duy nhất một $z \in G$ xác định bởi $g(f(x)) = z$

  \textbf{Hợp ánh xạ (tích ánh xạ)}: $(g \circ f)(x) = g(f(x))$. Và $(f \circ g)(x) = f(g(x))$. Nói chung $g \circ f \neq f \circ g$
  
  Điều kiện: $\mathrm{Im}(f) \subset$ miền xác định của $g$.
\end{dinhnghia}

\begin{tinhchat}[title={Tính chất hợp ánh xạ}]
  \begin{itemize}
    \item Tính kết hợp: $(h \circ g) \circ f = h \circ (g \circ f)$
    \item $f: A \to B$, $g: B \to C$ đơn ánh $\Rightarrow$ $g \circ f$ đơn ánh.
    \item $f: A \to B$, $g: B \to C$ toàn ánh $\Rightarrow$ $g \circ f$ toàn ánh.
    \item $f$, $g$ song ánh $\Rightarrow$ $g \circ f$ song ánh và $(g \circ f)^{-1} = f^{-1} \circ g^{-1}$.
  \end{itemize}
\end{tinhchat}

\begin{vidu}
  $f: \mathbb{R} \to \mathbb{R}^+$, $f(x) = e^x$ (song ánh). 
  
  $g: \mathbb{R}^+ \to \mathbb{R}$, $g(y) = \ln y$ (song ánh). 
  
  Khi đó $g \circ f = \mathrm{id}_\mathbb{R}$, $f \circ g = \mathrm{id}_{\mathbb{R}^+}$.

  Cho $f(x) = x^2$ và $g(x) = 2x + 5$. Khi đó: $(f \circ g)(x) = f(g(x)) = (2x + 5)^2 = 4x^2 + 20x + 25$ và $(g \circ f)(x) = 2x^2 + 5$
\end{vidu}

\begin{luuy}
    \begin{itemize}
        \item Nếu $f \circ g = \mathrm{id}_\mathbb{A}$ và $g \circ f = \mathrm{id}_{\mathbb{B}}$ thì $g$ là ánh xạ nghịch đảo của $f$
    \end{itemize}
\end{luuy}

\subsection{Thu hẹp và mở rộng}

\begin{dinhnghia}
  Cho $f: A \to B$, $C \subset A$. \textbf{Thu hẹp} của $f$ lên $C$: $f|_C: C \to B$, $f|_C(x) = f(x)$ với mọi $x \in C$.
\end{dinhnghia}

\begin{dinhnghia}
  $g: D \to B$ gọi là \textbf{mở rộng} của $f: A \to B$ nếu $A \subset D$ và $g|_A(x) = f(x)$ với mọi $x \in A$.
\end{dinhnghia}

\subsection{Lực lượng tập hợp}

\begin{dinhnghia}
  Hai tập $A$ và $B$ gọi là \textbf{đồng lực lượng} (cùng lực lượng) nếu tồn tại song ánh $f: A \to B$. Ký hiệu $|A| = |B|$ hoặc $A \sim B$.
\end{dinhnghia}

\begin{dinhnghia}
  Tập $A$ \textbf{hữu hạn} nếu $A = \varnothing$ hoặc $A \sim \{1, 2, \ldots, n\}$ với $n \in \mathbb{N}$. Khi đó $|A| = n$. 
  
  Tập \textbf{vô hạn}: không hữu hạn.
\end{dinhnghia}

\begin{dinhly}
  $\mathbb{N}$, $\mathbb{Z}$, $\mathbb{Q}$ đều đếm được (đồng lực lượng với $\mathbb{N}$). 
  
  $\mathbb{R}$ không đếm được (lực lượng lớn hơn $\mathbb{N}$).
\end{dinhly}

\begin{vidu}
  Song ánh $\mathbb{N} \to \mathbb{Z}$: $0 \mapsto 0$, $1 \mapsto 1$, $2 \mapsto -1$, $3 \mapsto 2$, $4 \mapsto -2$, $\ldots$ (xếp xen kẽ). Vậy $|\mathbb{Z}| = |\mathbb{N}|$.
\end{vidu}

%% ================================================================
%%  CHƯƠNG 5: CẤU TRÚC ĐẠI SỐ
%% ================================================================
\section{Cấu trúc đại số}

\subsection{Luật hợp thành trong trên một tập}

\begin{dinhnghia}
    Luật hợp thành trong trên tập $E$, hay phép toán trên $E$, là một quy luật khi tác động lên 2 phần tử $a, b \in E$ sẽ tạo ra 1 và chỉ 1 phần tử cũng $\in E$.

    Luật hợp thành trong trên tập $E$ là 1 ánh xạ từ $E \times E \to E$. Ký hiệu $*$.

    Ta có: $a, b \in E \to a * b \in E$
\end{dinhnghia}

\begin{tinhchat}[title={Tính chất của luật hợp thành trong}]
  \begin{itemize}
    \item \textbf{Giao hoán}: $a * b = b * a$.
    \item \textbf{Kết hợp}: $(a * b) * c = a * (b * c)$.
    \item \textbf{Phần tử trung hòa} $e$: $a * e = e * a = a$.
    \item \textbf{Phần tử đối xứng}: $\forall a \in E$, $\exists a' \in E$: $a * a' = a' * a = e$.
  \end{itemize}
\end{tinhchat}

\begin{dinhnghia}
    Một tập có trang bị 1 hay nhiều \textbf{luật hợp thành trong} với những tính chất xác định tạo thành 1 trong những đối tượng toán học gọi là \textbf{Cấu trúc đại số}
\end{dinhnghia}

\subsection{Nhóm}

\begin{dinhnghia}
  \textbf{Nhóm} $(G, *)$: tập $G$ với phép toán hai ngôi $*$ thỏa:
  \begin{enumerate}
    \item \textbf{Kết hợp}: $(a * b) * c = a * (b * c)$.
    \item \textbf{Phần tử trung hòa} $e$: $a * e = e * a = a$.
    \item \textbf{Phần tử đối xứng}: $\forall a \in G$, $\exists a' \in G$: $a * a' = a' * a = e$.
  \end{enumerate}
  Nếu thêm \textbf{giao hoán} $a * b = b * a$ thì gọi là \textbf{nhóm Abel} (nhóm giao hoán).
\end{dinhnghia}

\begin{tinhchat}[title={Tính chất của nhóm}]
  \begin{itemize}
    \item Phần tử trung hòa $e$ là duy nhất.
    \item Phần tử đối xứng $a'$ của $a$ là duy nhất (ký hiệu $a^{-1}$).
    \item Quy tắc giản ước: $a * x = a * y \Rightarrow x = y$.
    \item Phương trình $a * x = b$ có nghiệm duy nhất $x = a^{-1} * b$.
  \end{itemize}
\end{tinhchat}

\begin{vidu}
  $(\mathbb{Z}, +)$, $(\mathbb{Q} \setminus \{0\}, \cdot)$, nhóm các phép thế $S_n$.
\end{vidu}

\subsection{Vành}

\begin{dinhnghia}
  \textbf{Vành} $(R, +, \cdot)$: tập $R$ với hai phép toán $+$ và $\cdot$ thỏa:
  \begin{itemize}
    \item $(R, +)$ là nhóm Abel (nhóm giao hoán), thường phần tử trung hòa Luật cộng là $0$.
    \item Luật nhân ($\cdot$) có tính kết hợp.
    \item Luật nhân ($\cdot$) có tính phân phối 2 phía đối với luật cộng ($+$): $a \cdot (b + c) = a \cdot b + a \cdot c$, $(a + b) \cdot c = a \cdot c + b \cdot c$.
  \end{itemize}
  Nếu thêm Luật nhân ($\cdot$) có tính \textbf{giao hoán} $a \cdot b = b \cdot a$ thì gọi là Vành \textbf{giao hoán}.
  
  Vành \textbf{có đơn vị} nếu Luật nhân ($\cdot$) có phần tử trung hòa $1$: $1 \cdot a = a \cdot 1 = a$.
\end{dinhnghia}

\begin{vidu}
  $(\mathbb{Z}, +, \cdot)$, $(\mathbb{Q}, +, \cdot)$, $(\mathbb{R}, +, \cdot)$: vành giao hoán có đơn vị. 
  
  Tập các ma trận vuông cấp $n$: vành có đơn vị nhưng không giao hoán khi $n > 1$.
\end{vidu}

\begin{dinhnghia}
  \textbf{Vành nguyên} là vành $(R, +, \cdot)$ trong đó $a \cdot b = 0 \Rightarrow a = 0$ hoặc $b = 0$ (không có ước của 0).
\end{dinhnghia}

\subsection{Trường}

\begin{dinhnghia}
  \textbf{Trường} $(K, +, \cdot)$: tập $K$ với hai phép toán $+$ và $\cdot$ thỏa:
  \begin{itemize}
    \item Vành giao hoán có đơn vị
    \item Mọi phần tử khác $0$ (phần tử trung hòa của luật cộng ($+$)) đều có nghịch đảo đối với luật nhân ($\cdot$): $\forall a \in K$, $a \ne 0$, $\exists a^{-1}$: $a \cdot a^{-1} = 1$.
  \end{itemize}
\end{dinhnghia}

\begin{tinhchat}[title={Tính chất của trường}]
  \begin{itemize}
    \item Trường là vành nguyên.
    \item $K^* = K \setminus \{0\}$ là nhóm với phép nhân.
    \item Quy tắc giản ước: $a \ne 0$, $a \cdot x = a \cdot y \Rightarrow x = y$.
    \item Phương trình $a \cdot x = b$ ($a \ne 0$) có nghiệm duy nhất $x = a^{-1} \cdot b$.
  \end{itemize}
\end{tinhchat}

\begin{dinhnghia}
  \textbf{Trường con}: $F \subset K$ là trường con của $K$ nếu $F$ là trường với các phép toán cảm sinh từ $K$.
\end{dinhnghia}

\begin{vidu}
  $\mathbb{Q} \subset \mathbb{R} \subset \mathbb{C}$. $\mathbb{Q}$ là trường con nhỏ nhất của $\mathbb{R}$. $\mathbb{C} = \mathbb{R}(i)$ mở rộng bởi $i^2 = -1$.
\end{vidu}

\end{document}
